% =========================================================
% SUBIR ESTE CÓDIGO COMO CAPÍTULO DE PRÁCTICA EN LA MEMORIA
% =========================================================
\chapter{Cuaderno de Laboratorio — Práctica 2: Selección de Modelos y Regularización}

\section*{Objetivo}
Dominar el protocolo de \textbf{Validación Cruzada (K-Fold)} y comprender cómo la regularización
\textbf{Ridge (L2)} y \textbf{Lasso (L1)} previenen el sobreajuste y ayudan a la \textbf{selección de atributos}.

\section*{Material de partida}
Se parte del script \texttt{seleccion\_modelos\_viviendas.py} (o \texttt{.jl}), que incluye:
\begin{itemize}
    \item Un dataset sintético de \textbf{200 viviendas} con \textbf{10 variables predictoras}.
    \item Implementación de \textbf{validación cruzada} y del cálculo del \textbf{RMSE} (raíz del error cuadrático medio).
\end{itemize}

\vspace{0.4cm}
\noindent\textbf{Nota:} En este cuaderno se insertan capturas/gráficas generadas por el script (carpeta \texttt{outputs/}).
Ajusta los nombres de fichero a los que obtengas tú.

% =========================================================
% TAREA 1
% =========================================================
\section{Tarea 1 — Estabilidad de la Validación Cruzada (K-Fold)}

\subsection*{Configuración común}
Modelo usado: \underline{\hspace{6cm}} (regresión lineal sin regularización)\\
Métrica: \underline{\hspace{3cm}} (RMSE)\\
Semilla(s) / aleatoriedad: \underline{\hspace{6cm}}\\

\vspace{0.3cm}
\noindent\textbf{Hipótesis inicial (antes de ejecutar):}
¿Crees que el RMSE será estable al cambiar K? ¿Por qué?
\vspace{1.4cm}

\subsection{Experimento A — K = 2 (Validación simple): variabilidad entre ejecuciones}

\noindent Ejecuta el modelo \textbf{varias veces} con K=2 (por ejemplo 5--10) y registra el RMSE.

\vspace{0.2cm}
\begin{center}
\includegraphics[width=0.90\linewidth]{outputs/kfold_K2_rmse_runs.png}
\end{center}

\noindent\textbf{Registro de resultados (rellena):}
\begin{center}
\begin{tabular}{|c|c|}
\hline
Ejecución & RMSE \\
\hline
1 & \underline{\hspace{2.5cm}} \\
2 & \underline{\hspace{2.5cm}} \\
3 & \underline{\hspace{2.5cm}} \\
4 & \underline{\hspace{2.5cm}} \\
5 & \underline{\hspace{2.5cm}} \\
\hline
\end{tabular}
\end{center}

\vspace{0.2cm}
\noindent\textbf{Análisis:}
\begin{itemize}
    \item ¿Por qué el error varía tanto entre ejecuciones?
    \item ¿Es fiable esta métrica (con K=2) para tomar una decisión técnica?
\end{itemize}

\vspace{1.4cm}

\subsection{Experimento B — K = 10 (estándar): RMSE más estable}

\noindent Ejecuta el modelo con K=10 y anota el RMSE.

\vspace{0.2cm}
\begin{center}
\includegraphics[width=0.90\linewidth]{outputs/kfold_K10_rmse.png}
\end{center}

\noindent RMSE observado (K=10): \underline{\hspace{3cm}}

\vspace{0.4cm}
\noindent\textbf{Análisis:}
Describe si el resultado parece más estable que en K=2 y por qué.
\vspace{1.2cm}

\subsection{Experimento C — K = 100 (Leave-One-Out): coste vs. precisión}

\noindent Ejecuta el modelo con K=100 (LOOCV) y observa el RMSE y el tiempo.

\vspace{0.2cm}
\begin{center}
\includegraphics[width=0.90\linewidth]{outputs/kfold_K100_loocv_rmse.png}
\end{center}

\noindent RMSE observado (K=100): \underline{\hspace{3cm}} \quad
Tiempo aproximado: \underline{\hspace{3cm}}

\vspace{0.4cm}
\noindent\textbf{Análisis:}
Compara el coste computacional frente a K=10. ¿Merece la pena el tiempo extra para mejorar la precisión obtenida?
\vspace{1.2cm}

\subsection{Resumen comparativo (K = 2 vs 10 vs 100)}

\begin{center}
\begin{tabular}{|c|c|c|}
\hline
K & RMSE & Comentario de estabilidad/coste \\
\hline
2 & \underline{\hspace{2.5cm}} & \underline{\hspace{6cm}} \\
10 & \underline{\hspace{2.5cm}} & \underline{\hspace{6cm}} \\
100 & \underline{\hspace{2.5cm}} & \underline{\hspace{6cm}} \\
\hline
\end{tabular}
\end{center}

\vspace{0.3cm}
\noindent\textbf{Conclusión Tarea 1 (3--6 líneas):}
¿Cuál elegirías para este problema y por qué?
\vspace{1.6cm}

\newpage

% =========================================================
% TAREA 2
% =========================================================
\section{Tarea 2 — Regularización Ridge vs. Lasso (K = 10)}

\subsection*{Configuración}
Fijamos K = 10 y comparamos:
\begin{itemize}
    \item \textbf{Ridge Regression} (penalización L2)
    \item \textbf{Lasso Regression} (penalización L1)
\end{itemize}

\subsection{Ridge (L2) — $\lambda = 0.1,\, 10,\, 100$}

\noindent\textbf{Resultados:} captura el gráfico de barras de coeficientes para cada $\lambda$.

\vspace{0.2cm}
\begin{center}
\includegraphics[width=0.90\linewidth]{outputs/ridge_lambda_0.1_coeffs.png}
\end{center}
\noindent RMSE Ridge ($\lambda=0.1$): \underline{\hspace{3cm}}

\vspace{0.4cm}
\begin{center}
\includegraphics[width=0.90\linewidth]{outputs/ridge_lambda_10_coeffs.png}
\end{center}
\noindent RMSE Ridge ($\lambda=10$): \underline{\hspace{3cm}}

\vspace{0.4cm}
\begin{center}
\includegraphics[width=0.90\linewidth]{outputs/ridge_lambda_100_coeffs.png}
\end{center}
\noindent RMSE Ridge ($\lambda=100$): \underline{\hspace{3cm}}

\vspace{0.4cm}
\noindent\textbf{Análisis:}
¿Llegan a valer \textbf{cero} los coeficientes de las variables ruidosas (\texttt{Var\_4} a \texttt{Var\_10})?
Explica qué observas y por qué pasa (o no pasa).
\vspace{1.6cm}

\subsection{Lasso (L1) — $\lambda = 0.1,\, 10,\, 100$}

\noindent\textbf{Resultados:} captura el gráfico de barras de coeficientes para cada $\lambda$.

\vspace{0.2cm}
\begin{center}
\includegraphics[width=0.90\linewidth]{outputs/lasso_lambda_0.1_coeffs.png}
\end{center}
\noindent RMSE Lasso ($\lambda=0.1$): \underline{\hspace{3cm}}

\vspace{0.4cm}
\begin{center}
\includegraphics[width=0.90\linewidth]{outputs/lasso_lambda_10_coeffs.png}
\end{center}
\noindent RMSE Lasso ($\lambda=10$): \underline{\hspace{3cm}}

\vspace{0.4cm}
\begin{center}
\includegraphics[width=0.90\linewidth]{outputs/lasso_lambda_100_coeffs.png}
\end{center}
\noindent RMSE Lasso ($\lambda=100$): \underline{\hspace{3cm}}

\vspace{0.4cm}
\noindent\textbf{Análisis clave:}
\begin{enumerate}
    \item Identifica para qué valor de $\lambda$ el modelo \textbf{Lasso} consigue \textbf{eliminar} (poner a cero)
    las 7 variables ruidosas, quedándose solo con las 3 importantes:
    \[
    \lambda^{*} = \underline{\hspace{2.5cm}}
    \]
    \item Variables con coeficiente distinto de cero (modelo Lasso con $\lambda^{*}$):\\
    \underline{\hspace{13cm}}
\end{enumerate}

\vspace{1.0cm}

\subsection{El dilema de la complejidad — $\lambda$ demasiado alto (ej. 1000)}

\noindent Prueba un valor extremo, por ejemplo $\lambda = 1000$, y registra qué ocurre.

\vspace{0.2cm}
\begin{center}
\includegraphics[width=0.90\linewidth]{outputs/lasso_lambda_1000_coeffs.png}
\end{center}

\noindent RMSE ($\lambda=1000$): \underline{\hspace{3cm}}

\vspace{0.4cm}
\noindent\textbf{Análisis:}
Explica qué ocurre con el RMSE si aplicas un $\lambda$ demasiado alto.
¿Estamos ante un caso de \textbf{sesgo alto} o \textbf{varianza alta}? Justifica.
\vspace{1.8cm}

\newpage

% =========================================================
% RETO
% =========================================================
\section{Reto — El “Cazador de Variables”}

\noindent\textbf{Contexto:} Imagina que cada variable incluida en tu modelo tiene un \textbf{coste de mantenimiento}.
Tu objetivo es encontrar el modelo \textbf{más sencillo posible} (menos variables con coeficientes $\neq 0$)
que mantenga un RMSE inferior a un umbral:

\[
\text{RMSE} < \underline{\hspace{2.5cm}} \quad (\text{umbral del profesor})
\]

\vspace{0.3cm}
\noindent\textbf{Mejor solución encontrada (rellena):}
\[
\lambda = \underline{\hspace{2.5cm}}
\quad\quad
\text{\#Variables activas} = \underline{\hspace{2.5cm}}
\quad\quad
\text{RMSE} = \underline{\hspace{2.5cm}}
\]

\vspace{0.3cm}
\noindent Variables seleccionadas (coeficiente $\neq 0$):
\underline{\hspace{13cm}}

\vspace{0.5cm}
\begin{center}
\includegraphics[width=0.90\linewidth]{outputs/reto_best_model_coeffs.png}
\end{center}

\vspace{0.3cm}
\noindent\textbf{Justificación técnica (6--10 líneas):}
¿Por qué este modelo es el mejor compromiso entre simplicidad (coste) y rendimiento (RMSE)?
\vspace{2.2cm}

% =========================================================
% CONCLUSIONES
% =========================================================
\section{Conclusiones finales}
Resume en 8--12 líneas lo aprendido sobre:
\begin{itemize}
    \item Estabilidad de K-Fold al variar K (K=2 vs K=10 vs LOOCV).
    \item Diferencias prácticas entre Ridge y Lasso (coeficientes, selección de variables).
    \item Relación entre $\lambda$ grande y el sesgo/varianza.
\end{itemize}

\vspace{3.0cm}
